\paragraph{\trHomotopy{}与\trIsomorphism。}\label{sub:homotopy-isomorphism}
设 $\scr{A}$ 是\trStrictlyUnital{}\trACategory，$\scr{C}$ 是\trNonunital{}的，而
$\scr{F}_0,\scr{F}_1\colon\scr{C}\to\scr{A}$ 是在\trObject{}上作用相同的%
\trNonunital{}\trAFunctor。设 $T$ 是从 $\scr{F}_0$ 到 $\scr{F}_1$ 的一个\trDegree{}为
零的预\trNaturalTransformation，满足 $T^0=0$。用 $S$ 表示修改后的预%
\trNaturalTransformation，它由 $S^0=e_{\scr{F}_0(X)}$（对一切 $X$）、$S^d=T^d$
（$d>0$）定义。一个直截了当的计算表明：$T$ 是 $\scr{F}_0$ 与 $\scr{F}_1$ 之间的%
\trHomotopy，当且仅当 $S$ 是一个真正的\trNaturalTransformation。在这一语境下，
\trHomotopy{}之间的\trComposition{}表达式\eqref{eq:transitivity-homotopy} 这个原本
颇为神秘的东西，可以看出无非就是相应\trNaturalTransformation{}的通常\trComposition；
对\eqref{eq:symmetry-maps-phi-psi} 亦然。

\begingroup\tolerance=9999\emergencystretch=3em\hbadness=10000
利用\cref{thm:right-composition-iso} 以及 \trAinfty-\trFunctor\trCategory{}的%
\trStrictlyUnital{}性，可知两个\trHomotopic{}的\trFunctor{}作为\penalty0%
$H^0(\mathit{nu}\text{-}\mathit{fun}(\scr{C},\scr{A}))$ 中的\trObject{}是%
\trIsomorphic{}的。我们将会在\trCunital{}的语境下需要与此类似的结果：\par\endgroup
\vkthmsync{5}
\begin{lemma}\label{thm:homotopic-functors-isomorphic}
设 $\scr{A}$ 是\trCunital{}\trACategory，$\scr{C}$ 是\trNonunital{}的。若两个%
\trNonunital{}\trAFunctor{} $\scr{F}_0,\scr{F}_1\colon\scr{C}\to\scr{A}$ 是%
\trHomotopic{}的，则它们作为 $H^0(\mathit{nu}\text{-}\mathit{fun}(\scr{C},\scr{A}))$
中的\trObject{}也是\trIsomorphic{}的。
\end{lemma}
\begin{proof}
我们沿用\cref{thm:functor-category-cunital} 证明中的记号。由于左%
\trComposition{}保持\trHomotopy，诸 $\tilde{\scr{F}}_k=\Phi\circ\scr{F}_k$ 是%
\trHomotopic{}的，于是依上面的观察，它们作为 $H^0(\tilde{\scr{Q}})$ 中的\trObject{}%
是\trIsomorphic{}的。另一方面，$\scr{L}_{\Phi}$ 由\cref{thm:left-composition-cff}
知是\trCohomology\trFullyFaithful{}的，故诸 $\scr{F}_k$ 自身必是\trIsomorphic{}的。
\end{proof}
